\section{갈루아 이론으로 가는 다리}
\label{sec:polynomial-galois-bridge}

다항식환을 배우는 목적은 단순히 인수분해 계산을 잘하기 위해서가 아니다. 방정식
\[
  f(x)=0
\]
을 구조적으로 연구하려면 다음 질문에 답해야 한다.

\begin{enumerate}
  \item $f$는 계수체 $K$ 위에서 더 작은 차수의 식으로 분해되는가?
  \item 분해되지 않는다면, $f$의 근을 포함하는 체를 어떻게 만들 수 있는가?
  \item 서로 다른 근을 모두 포함하는 최소의 체는 무엇인가?
  \item 그 체의 $K$-자동동형들은 근들을 어떻게 바꾸는가?
\end{enumerate}

첫 질문은 기약성의 문제이다. 두 번째 질문은 몫체
\[
  K[X]/(f)
\]
가 답한다. 세 번째 질문은 다음 권의 분해체 이론으로 이어지고, 네 번째 질문이 갈루아군을 낳는다.

\begin{keyidea}
기약다항식 $f\in K[X]$의 한 근 $\alpha$를 어떤 확대체에서 택하면 $f$는 $\alpha$의 최소다항식이고
\[
 K(\alpha)\cong K[X]/(f)
\]
이다. 이후에는 $f$의 나머지 근을 첨가하고, 그 근들을 움직이면서 $K$를 고정하는 자동동형을 연구한다.
\end{keyidea}

\begin{example}
$f=X^3-2\in\Q[X]$는 아이젠슈타인 판정법으로 기약이다. 따라서
\[
  L=\Q[X]/(X^3-2)
\]
는 체이고, $\alpha=X+(X^3-2)$라 두면 $\alpha^3=2$이다. 모든 원소는
\[
  a+b\alpha+c\alpha^2\qquad(a,b,c\in\Q)
\]
로 유일하게 표현된다. 이 추상적 몫체를 실수의 세제곱근 $\sqrt[3]{2}$에 대입하여 $\Q(\sqrt[3]{2})$와 동형으로 볼 수 있다. 그러나 이 체는 $f$의 세 복소근을 모두 포함하지 않는다. 나머지 복소근을 포함하려면 더 큰 분해체가 필요하다.
\end{example}

\clearpage
\section{장 요약}

\begin{chapterreview}
\begin{enumerate}
  \item 다항식은 다항함수가 아니라 유한 지지 계수열로 정의한다.
  \item 정역 $R$에서는 $\deg(fg)=\deg f+\deg g$이고 $R[X]$도 정역이다.
  \item 나누는 다항식의 최고차항계수가 단위원이면 나눗셈 알고리즘이 성립한다.
  \item $K[X]$는 유클리드 정역이므로 최대공약수와 베주 항등식을 계산할 수 있고 UFD이다.
  \item $f(a)=0$일 필요충분조건은 $(X-a)\mid f$이다.
  \item 영이 아닌 $f$의 중근은 $f$와 $f'$의 공통근이며, 제곱인수가 없을 조건은 $\gcd(f,f')=1$이다.
  \item 차수 $2$ 또는 $3$에서는 근의 존재가 가약성과 동치이다.
  \item 비상수 $f\in K[X]$가 기약일 필요충분조건은 $K[X]/(f)$가 체인 것이다.
  \item 양의 차수인 원시다항식의 $R[X]$에서의 기약성과 $Q(R)[X]$에서의 기약성은 동치이다.
  \item 법 $p$ 축소와 아이젠슈타인은 강력한 기약성 충분조건이다.
  \item $R$이 UFD이면 여러 변수 다항식환도 UFD이지만 일반적으로 PID는 아니다.
\end{enumerate}
\end{chapterreview}

\section{빠른 복습 카드}

\begin{description}[style=nextline]
  \item[Q] 다항식과 다항함수가 항상 같은 대상이 아닌 이유는?\\
  \textbf{A} 유한체에서는 서로 다른 다항식이 같은 함수를 정의할 수 있기 때문이다.

  \item[Q] $K[X]$에서 나눗셈 알고리즘이 항상 가능한 이유는?\\
  \textbf{A} 영이 아닌 최고차항계수가 체의 단위원이기 때문이다.

  \item[Q] $m\ne0$일 때 $\overline f\in K[X]/(m)$가 단위원일 조건은?\\
  \textbf{A} $\gcd(f,m)=1$.

  \item[Q] 차수 $2$ 또는 $3$인 다항식의 기약성 검사법은?\\
  \textbf{A} 계수체 안에 근이 있는지 검사한다.

  \item[Q] 비상수 $f\in K[X]$에 대해 $K[X]/(f)$가 체일 조건은?\\
  \textbf{A} $f$가 $K[X]$에서 기약일 것.

  \item[Q] 아이젠슈타인 조건의 핵심은?\\
  \textbf{A} 최고차항계수는 $p$로 나누어지지 않고, 나머지는 모두 나누어지며, 상수항은 $p^2$으로 나누어지지 않는 것.
\end{description}
